\documentclass[10pt,a4paper]{amsart}
\usepackage[margin=19mm]{geometry}
\usepackage{amsmath,amssymb,mathtools}
\usepackage[scr=rsfs,cal=euler]{mathalfa}
\usepackage{xfrac}
\usepackage[unicode,hidelinks,bookmarks=false]{hyperref}
\newcommand{\ie}{i.e.\ }
\newcommand{\cf}{cf.\ }
\newcommand{\resp}{respectively\ }
\newcommand{\Subsection}{\subsection}
\providecommand{\nobreakdash}{\nobreak\hbox{-}}
\setlength{\emergencystretch}{2em}
\multlinegap0pt
\usepackage{bbold,mathtools,bm}
\def\-{\raisebox{.75pt}{-}}
\usepackage{enumitem}
\usepackage{tikz}
\usetikzlibrary{cd}
\usetikzlibrary{arrows}
\numberwithin{table}{section}
\numberwithin{equation}{section}
\newtheorem{defn}{Definition}[section]
\newtheorem{exmp}{Example}[section]
\newtheorem{rmk}{Remark}[section]
\newtheorem{lem}{Lemma}[section]
\newtheorem{prop}{Proposition}[section]
\newtheorem{thm}{Theorem}[section]
\newtheorem*{thm*}{Theorem}
\newtheorem{cor}{Corollary}[section]
\newtheorem{obs}{Observation}[section]
\DeclareMathOperator{\BV}{BV}
\newcommand{\tr}{\mathrm{tr}}
\newcommand{\bbA}{\mathbb{A}}
\newcommand{\bbB}{\mathbb{B}}
\newcommand{\fg}{\mathfrak{g}}
\newcommand{\fh}{\mathfrak{h}}
\newcommand{\fk}{\mathfrak{k}}
\newcommand{\fl}{\mathfrak{l}}
\newcommand{\fw}{\mathfrak{w}}
\newcommand{\sD}{\mathcal{D}}
\newcommand{\sE}{\mathcal{E}}
\newcommand{\sF}{\mathcal{F}}
\newcommand{\sG}{\mathcal{G}}
\newcommand{\sH}{\mathcal{H}}
\newcommand{\sL}{\mathcal{L}}
\newcommand{\sM}{\mathcal{M}}
\newcommand{\K}{\mathbb{K}}
\newcommand{\R}{\mathbb{R}}
\newcommand{\C}{\mathbb{C}}
\newcommand{\Hq}{\mathbb{H}}
\newcommand{\Z}{\mathbb{Z}}
\newcommand{\N}{\mathbb{N}}
\newcommand{\Q}{\mathbb{Q}}
\newcommand{\ut}{\mathfrak{u}}
\newcommand{\U}{\mathrm{U}}
\newcommand{\Ad}{\mathrm{Ad}}
\newcommand{\diag}{\mathrm{diag}}
\newcommand{\GL}{\mathrm{GL}}
\newcommand{\sdet}{\mathrm{sdet}\,}
\newcommand{\rad}{\mathrm{rad}}
\newcommand{\Tor}{\mathrm{Tor}}
\newcommand{\sgn}{\mathrm{sgn}\,}
\newcommand{\Aut}{\mathrm{Aut}}
\newcommand{\Der}{\mathrm{Der}}
\newcommand{\vol}{\mathrm{vol}}
\newcommand{\covol}{\mathrm{covol}}
\newcommand{\dive}{\mathrm{div}}
\newcommand{\Lie}{\mathrm{Lie}}
\newcommand{\Ber}{\mathrm{Ber}}
\usepackage{cleveref}
\usepackage{amsrefs}
\renewcommand{\eprint}[1]{\href{http://www.arxiv.org/abs/#1}{\texttt{arXiv:#1}}}
\begin{document}
\title[Higher abelian gauge theory in BV formalism]{Higher abelian gauge theory in BV formalism}
\author{Shuhan Jiang}
\date{}
\maketitle
\noindent\emph{Source and licence.} Higher abelian gauge theory in BV formalism. Version arXiv:2609.17518v1 (CC BY 4.0). This material is distributed under the Creative Commons Attribution 4.0 International licence, http://creativecommons.org/licenses/by/4.0/, as recorded in the arXiv OAI licence metadata for this exact version and shown on the abstract page https://arxiv.org/abs/2609.17518v1. Full rights notice: licenses/arxiv-2609.17518-CC-BY-4.0.txt.\par\smallskip
\noindent\textit{Editorial scope.} Counted unit: the original \texttt{thm} environment \texttt{thm:L-independence} at source lines 458--460 together with its complete proof at lines 462--498; the delivered document keeps source lines 148--499 so that every referenced label and every citation resolves inside it. Native editable source is the paper's own concatenated TeX; the AMS wrapper is editorial formatting only. No publisher class, upstream script or unrelated artwork is redistributed.
\section{Abelian BV groups}

In this section, we introduce an abstract BV structure on a collection of abelian Lie groups, modeled on the global zero modes of higher abelian gauge theories.

\subsection{Preliminaries}

Let $V = \bigoplus_{k\in \Z} V^k$ be a finite-dimensional graded vector space.

\begin{defn}
	A \emph{$(-1)$-shifted symplectic pairing} on $V$ is a nondegenerate graded skew-symmetric pairing
	\[
	\omega \colon V \times V \longrightarrow \R[-1].
	\]
	
	Graded skew-symmetry means that $\omega(V^k,V^\ell)=0$ unless $k+\ell=1$, and
	\[
	\omega(x,y)=-(-1)^{(k-1)(\ell-1)}\omega(y,x).
	\]
	for $x \in V^k$ and $y \in V^\ell$. 
	
	Nondegeneracy means that the map
	\[
	\omega^\flat \colon V \xlongrightarrow{\sim} V^\vee[-1],
	\qquad x\longmapsto\omega(x,-),
	\]
	is an isomorphism.
\end{defn}

For a graded subspace $W\subset V$, let
\[
W^{\perp_\omega} \coloneq \{x\in V\mid \omega(x,W)=0\}.
\]
We call $W$ \emph{isotropic}, \emph{coisotropic}, or \emph{Lagrangian} if, respectively,
\[
W\subset W^{\perp_\omega},
\qquad
W^{\perp_\omega}\subset W,
\qquad
W=W^{\perp_\omega}.
\]

\begin{defn}
	A (graded) lattice $\Gamma \subset V$ is called \emph{integral} if
	\[
	\omega(\Gamma,\Gamma)\subset\Z.
	\]
	It is called \emph{full} if
	\[
	\Gamma_\R \coloneqq  \Gamma\otimes_{\Z}\R \cong V.
	\]
    Its \emph{radical} is
	\[
	\rad(\Gamma)\coloneq\{\gamma\in\Gamma\mid \omega(\gamma,\Gamma)=0\}.
	\]
\end{defn}

An integral lattice is \emph{coisotropic} if $\Gamma_\R$ is coisotropic. In this case, $\rad(\Gamma)$ is a full lattice in $\Gamma_\R^{\perp_\omega}$, and $\omega$ induces a nondegenerate integral pairing $\underline\omega$ on
\[
\underline\Gamma \coloneq \Gamma/\rad(\Gamma).
\]
We call a coisotropic lattice $\Gamma$ \emph{reduced unimodular} if
\[
\underline \omega^\flat:
\underline\Gamma\xlongrightarrow{\sim}
\underline\Gamma^\vee[-1]
\]
is an isomorphism.

Let $A$ be a finitely generated abelian group. 

\begin{defn}
A complex \emph{quadratic weight} on $A$ is a map
\[
w:A\longrightarrow\C/i\Z
\]
such that
\begin{itemize}
	\item $w(x+y)-w(x)-w(y)$ is biadditive;
	\item $\Re\bigl(w(x)\bigr)=\Re\bigl(w(-x)\bigr)$.
\end{itemize}
\end{defn}

Because the period group $i \Z \subset \C$ is purely imaginary, the function
\[
q(x)\coloneq\Re\bigl(w(x)\bigr)
\]
is a well-defined real quadratic weight on $A$.

\begin{lem}\label{lem:real-quadratic}
The function $q$ vanishes on $\Tor(A)$ and descends uniquely to a real quadratic weight
\[
\bar q:A_{\mathrm{free}}\coloneq A/\Tor(A)\longrightarrow\R
\]
satisfying $\bar q(nx)=n^2\bar q(x)$ for every $n\in\Z$.
\end{lem}

\begin{proof}
Let $b(x,y)\coloneqq q(x+y)-q(x)-q(y)$. Since $q(-x)=q(x)$, the identity $0=q(x-x)=2q(x)-b(x,x)$ gives $b(x,x)=2q(x)$. Since $2 \in \R$ is invertible, $q(nx)= b(nx,nx)/2=n^2q(x)$. If $t$ is torsion and $nt=0$, then $0=q(nt)=n^2q(t)$, so $q(t)=0$. Moreover $n b(x,t)=b(x,nt)=0$, and therefore $b(x,t)=0$. Thus $q(x+t)=q(x)$, proving that $q$ descends to $A_{\mathrm{free}}$.
\end{proof}

\subsection{BV structures}

Let $H=\{H^k\}_{k\in\Z}$ be a collection of abelian Lie groups, trivial in all but finitely many degrees, such that each Lie algebra $\fh^k$ of $H^k$ is finite-dimensional and each component group $\pi_0(H^k)$ is finitely generated.

Write
\[
\fh \coloneq \bigoplus_{k\in\Z}\fh^k,
\qquad
\Gamma^k \coloneq \ker\bigl(\exp:\fh^k\to H^k\bigr),
\qquad
\Gamma \coloneq \bigoplus_k\Gamma^k.
\]
For every $k$ there is a canonical exact sequence
\begin{equation}\label{eq:exp-sequence}
	0\longrightarrow \Gamma^k
	\longrightarrow \fh^k
	\xlongrightarrow{\exp} H^k
	\longrightarrow \pi_0(H^k)
	\longrightarrow 0.
\end{equation}
Let $G_0$ denote the identity component of an abelian Lie group $G$. Then the real span $\Gamma^k_\R$ of $\Gamma^k$ can be identified with the Lie algebra of the maximal compact subgroup of $H_0^k$. In particular, $H_0^k$ is compact if and only if $\Gamma^k$ is a full lattice in $\fh^k$.

We set
\[
H^{\mathrm{even}} \coloneq \prod_{k\in2\Z}H^k,
\qquad
H^{\mathrm{odd}} \coloneq \prod_{k\in2\Z+1}H^k,
\]
and use the same convention for $\fh$, $\Gamma$, and subgroups of $H$.


\begin{defn}\label{def:abelian-BV-group}
A \emph{BV structure} on $H$ is a triple $(\omega,\kappa,w)$, where:
\begin{itemize}[leftmargin=2em]
\item $\omega$ is a $(-1)$-shifted symplectic pairing on $\fh$;
\item $\kappa = \bigoplus \kappa^k$ is a direct sum of inner products $\kappa^k$ on $\fh^k$ for which
\[
\omega^\flat:(\fh,\kappa)\xlongrightarrow{\sim}(\fh^\vee[-1],\kappa^\vee[-1])
\]
is an isometry;
\item $w$ is a complex quadratic weight on $\pi_0(H^0)$. 
\end{itemize}
We call $H$, equipped with such a BV structure, an \emph{abelian BV group}.
\end{defn}

We call an abelian BV group $H$ \emph{integral} if its exponential lattice $\Gamma$ is integral with respect to the $(-1)$-shifted symplectic pairing $\omega$. An integral abelian BV group is called \emph{integrable} if $\Gamma$ is, in addition, coisotropic.

\begin{defn}
	A \emph{Lagrangian} of $H$ is a family of closed subgroups
	\[
	L=\{L^k\subset H^k\}_{k\in\Z}
	\]
	such that the total Lie algebra $\fl\subset\fh$ is Lagrangian and the induced map
	\[
	\pi_0(L^k)\longrightarrow\pi_0(H^k)
	\]
	is injective for all $k\in\Z$.
\end{defn}

A Lagrangian $L$ is called \emph{admissible} if each $L_0^k$ is compact, $\pi_0(L^k)$ is finite for $k\neq 0$, and
\[
\Re(w)\big|_{\pi_0(L^0)}
\]
descends to a positive-definite quadratic form on $\pi_0(L^0)/\Tor(\pi_0(L^0))$. 

\begin{lem}
	Every integrable abelian BV group $H$ has an admissible Lagrangian.
\end{lem}

\begin{proof}
	Choose a Lagrangian subspace $\underline{\fl}$ of the reduced symplectic space $\Gamma_\R/\rad(\Gamma)_\R$ rational with respect to $\underline\Gamma = \Gamma/\rad(\Gamma)$, and let $\fl \subset\Gamma_\R$ be its inverse
	image. Then $\fl$ is Lagrangian in $\fh$ and $\Gamma\cap\fl$ is a full lattice
	in $\fl$. Hence
	\[
	L \coloneq \exp(\fl)\subset H_0
	\]
	is a compact connected Lagrangian. In particular, $\pi_0(L)=0$, so $L$ is admissible.
\end{proof}

Henceforth, we assume that $H$ is integrable.


\begin{lem}\label{lem:integrable-admissible}
	A Lagrangian $L$ is admissible if and only if
	\[
	\int_{L^{\mathrm{even}}} \left|\exp(-2\pi w)\right|\,\mu_{\mathrm{even}}<\infty,
	\qquad
	\int_{L^{\mathrm{odd}}}1\ \mu_{\mathrm{odd}}<\infty,
	\]
	where $\mu_{\mathrm{even}}$ and $\mu_{\mathrm{odd}}$ denote the positive densities on $L^{\mathrm{even}}$ and $L^{\mathrm{odd}}$ induced by $\kappa$.
\end{lem}

\begin{proof}
	Since $|\exp(-2\pi w)|=\exp(-2\pi\Re(w))$ is constant on connected components and $\mu_{\mathrm{even}}$ is translation-invariant,
	\[
	\int_{L^{\mathrm{even}}} \left|\exp(-2\pi w)\right|\,\mu_{\mathrm{even}}
	=
	\vol(L^{\mathrm{even}}_0)
	|\pi_0(L^{\mathrm{even}})/\pi_0(L^0)|
	\sum_{c\in\pi_0(L^0)}\exp(-2\pi\Re(w(c))).
	\]
	The volume factor is finite precisely when $L^{\mathrm{even}}_0$ is compact. The second factor is finite precisely when $\pi_0(L^{\mathrm{even}})/\pi_0(L^0)$ is finite. By Lemma~\ref{lem:real-quadratic}, the series is a finite torsion factor times a Gaussian sum over the free lattice. It converges if and only if the induced real quadratic form is positive definite.
		
	Likewise,
	\[
	\int_{L^{\mathrm{odd}}}1\ \mu_{\mathrm{odd}} = |\pi_0(L^{\mathrm{odd}})|
	\vol(L^{\mathrm{odd}}_0),
	\]
	which is finite precisely when the identity component $L^{\mathrm{odd}}_0$ is compact and the component group $\pi_0(L^{\mathrm{odd}})$ is finite.
\end{proof}

An admissible Lagrangian is called  \emph{maximal} if it is not properly contained in another admissible Lagrangian.

\begin{defn}
For a maximal admissible Lagrangian $L$, define
\begin{equation}\label{eq:BV-partition}
Z_H(L)\coloneq
\frac{\displaystyle\int_{L^{\mathrm{even}}}\exp(- 2\pi w)\,\mu_{\mathrm{even}}}
{\displaystyle\int_{L^{\mathrm{odd}}}1\ \mu_{\mathrm{odd}}}.
\end{equation}
We call $Z_H(L)$ the \emph{zero-mode partition function} of $H$ (with respect to $L$).
\end{defn}

The definition is motivated by two elementary features of the field-theoretic partition function. First, it is modeled on finite-dimensional BV integration over a Lagrangian. Second, partition functions are multiplicative under disjoint unions. Since the zero modes of a disjoint union form the product of the corresponding abelian BV groups, the quotient of the even and odd integrals in \eqref{eq:BV-partition} has the expected multiplicativity. 

For convenience, we set
\[
T_k\coloneq\bigl|\Tor (\pi_0(H^k))\bigr|,
\qquad
T(H)\coloneq\prod_{k\ne0}T_k^{(-1)^k},
\]
\[
\Theta(L)\coloneq\sum_{c\in\pi_0(L^0)}\exp(-2 \pi w(c))
\qquad
\vol(L_0)\coloneq\prod_k \vol(L_0^k)^{(-1)^k},
\]
with the convention $\vol(\{0\})=1$.

\begin{prop}\label{prop:partition-factorization}
For every maximal admissible Lagrangian $L$,
\begin{equation}\label{eq:partition-factorization}
Z_H(L)=\vol(L_0)\, T(H)\, \Theta(L).
\end{equation}
\end{prop}

\begin{proof}
Enlarging a Lagrangian by discrete components does not change its Lie algebra. Maximality therefore forces
\[
\pi_0(L^k)=
\Tor(\pi_0(H^k)), \qquad k\ne0.
\]
The rest follows from the proof of Lemma \ref{lem:integrable-admissible}.
\end{proof}

For a Euclidean vector space $(V,\kappa)$ and a full lattice $\Lambda\subset V$, write
\[
{\det}_\Lambda \kappa\coloneq\det\bigl(\kappa(e_i,e_j)\bigr),
\qquad
\covol(\Lambda)\coloneq\sqrt{{\det}_\Lambda \kappa},
\]
where $(e_i)$ is any basis of $\Lambda$. Then
\[
\vol(V/\Lambda) = \covol(\Lambda).
\]
For graded lattices, write $\covol(\Lambda)\coloneqq\prod_k\covol(\Lambda^k)^{(-1)^k}$.


\begin{lem}\label{lem:unimodular-volume}
Suppose that $\Gamma$ is a full unimodular lattice in $\fh$, and let $\fl\subset\fh$ be a Lagrangian such that 
\[
\Gamma_L\coloneq\Gamma\cap\fl
\]
is full in $\fl$. Then
\[
\vol(L_0)=\bigl({\det}_\Gamma \kappa \bigr)^{1/4},
\qquad
{\det}_\Gamma \kappa \coloneq\prod_k\bigl({\det}_{\Gamma^k}\kappa^k\bigr)^{(-1)^k}.
\]
In particular, the graded volume $\vol(L_0)$ is independent of the embedding $L_0\hookrightarrow H_0$.
\end{lem}

\begin{proof}
Unimodularity and the Lagrangian condition give an exact sequence of graded lattices
\[
0\longrightarrow\Gamma_L\longrightarrow\Gamma
\longrightarrow\Gamma_L^\vee[-1]\longrightarrow0.
\]
Taking Berezinian lines yields
\[
\Ber_\Z(\Gamma)\cong
\Ber_\Z(\Gamma_L)\otimes\Ber_\Z(\Gamma_L^\vee[-1])
\cong\Ber_\Z(\Gamma_L)^{\otimes2}.
\]
Compatibility of $\kappa$ and $\omega$ makes this identification an isometry. Taking norms of the integral Berezinian generators $1_\Gamma \in \Ber_\Z(\Gamma)$ and $1_{\Gamma_L} \in \Ber_\Z(\Gamma_L)$ then gives
\[
({\det}_{\Gamma}\kappa)^{1/2}
=\lVert 1_{\Gamma}\rVert
=\lVert 1_{\Gamma_L}\rVert^2
={\det}_{\Gamma_L}\kappa|_{\Gamma_L}
=\vol_{\kappa}(L_0)^2.
\]
and the result follows.
\end{proof}

We call $H$ \emph{nice} if it satisfies the following additional conditions:
\begin{itemize}
	\item $\Gamma$ is reduced unimodular;
	\item $q=\Re(w)$ descends to a positive-definite quadratic form on $\pi_0(H^0)/\Tor(\pi_0(H^0))$.
\end{itemize}


\begin{thm}\label{thm:L-independence}
The zero-mode partition function $Z_H(L)$ of a nice $H$ is independent of the maximal admissible Lagrangian $L$.
\end{thm}

\begin{proof}
By Proposition~\ref{prop:partition-factorization}, $Z_H(L)=\vol(L_0)\, T(H)\, \Theta(L)$. The torsion factor $T(H)$ is independent of $L$. Since $H$ is nice, $q=\Re(w)$ is positive definite on the real span of the free part of $\pi_0(H^0)$. Maximality therefore forces
\[
\pi_0(L^0)=\pi_0(H^0).
\]
Hence $\Theta(L)$ is also independent of $L$. 

It remains only to prove that $\vol(L_0)$ is independent of $L$. 
Since $L_0$ is compact, its Lie algebra satisfies $\fl\subset\Gamma_\R$. As $\Gamma_\R^{\perp_\omega} \subset \Gamma_\R$ and $\fl$ is Lagrangian,
\[
\Gamma_\R^{\perp_\omega} \subset \fl^{\perp_\omega}=\fl.
\]
Consequently $\Gamma\cap \Gamma_\R^{\perp_\omega} \subset\Gamma_L=\Gamma\cap\fl$. Coisotropic reduction then gives an inclusion of lattices
\[
\underline\Gamma_L \coloneq \Gamma_L/(\Gamma\cap \Gamma_\R^{\perp_\omega} )
\subset
\underline\Gamma = \Gamma/\rad(\Gamma).
\]
The latter is unimodular by assumption. Equip $\Gamma_\R^{\perp_\omega}$ and $\underline{\Gamma}_\R = \Gamma_\R/\Gamma_\R^{\perp_\omega}$ with the inner products induced by $\kappa$. Graded covolumes are multiplicative for the resulting short exact sequence of lattices
\[
0 \longrightarrow \Gamma\cap \Gamma_\R^{\perp_\omega} \longrightarrow \Gamma_L \longrightarrow \underline\Gamma_L \longrightarrow 0,
\]
which is full in the short exact sequence of Euclidean vector spaces
\[
0 \longrightarrow \Gamma_\R^{\perp_\omega} \longrightarrow (\Gamma_L)_\R = \fl \longrightarrow (\underline \Gamma_L)_\R \longrightarrow 0.
\]
It follows that
\[
\vol(L_0)
=
\covol(\Gamma_L)
=
\covol(\Gamma \cap \Gamma_\R^{\perp_\omega})\,
\covol(\underline\Gamma_L).
\]
The first factor is manifestly independent of $L$, while Lemma~\ref{lem:unimodular-volume} implies that the second factor is also independent of $L$.
\end{proof}

\end{document}
